% Capítulo 1: Introducción y Motivación
% Cubre las secciones 1 y 2 de la memoria

\section{Introducción}

Las redes neuronales profundas dominan actualmente el estado del arte en tareas de visión por
computador, procesamiento del lenguaje natural y reconocimiento de patrones. Sin embargo, el
paradigma dominante de entrenamiento (procesar todos los datos de una vez con acceso
irrestricto) resulta inapropiado en muchos escenarios reales. Los datos llegan de forma
secuencial, los recursos computacionales son limitados, y las restricciones de privacidad o
almacenamiento impiden conservar todos los datos históricos. En este contexto surge el campo del
\textbf{aprendizaje continuo} (\textit{continual learning}), cuyo objetivo es que un modelo aprenda
de nuevos datos sin olvidar los conocimientos previos.

El obstáculo principal de este campo es el \textbf{olvido catastrófico}
(\textit{catastrophic forgetting}), descrito por primera vez de forma sistemática por McCloskey y
Cohen en 1989~\cite{ref1}. Cuando una red neuronal es re-entrenada sobre nuevas muestras mediante
descenso de gradiente, la actualización de sus pesos interfiere destructivamente con las
representaciones aprendidas anteriormente, de forma mucho más severa que el olvido observado en
sistemas biológicos.

La mayoría de las soluciones propuestas en la literatura (regularización, expansión de
arquitectura, o replay de datos generados) introducen compromisos entre plasticidad (capacidad de
aprender cosas nuevas) y estabilidad (capacidad de retener las anteriores), y frecuentemente
requieren almacenar ejemplos pasados o aumentar la complejidad del modelo con cada nueva tarea.

Este trabajo evalúa Qsimov, una librería que implementa un enfoque alternativo basado en una
propiedad matemática de las redes con activación ReLU: para cualquier entrada, la función que
computa la red es \textbf{lineal} en la región del espacio que contiene a esa entrada. Cada entrada
traza un <<camino>> único a través de la red, determinado por qué neuronas están activas ($> 0$) y
cuáles no. Qsimov separa la estructura de esos caminos, aprendida durante el pre-entrenamiento,
de los pesos asociados a cada camino, que se resuelven analíticamente mediante un sistema de
ecuaciones lineales. La acumulación incremental de ecuaciones a lo largo del tiempo garantiza que el
sistema <<recuerde>> todas las rondas de entrenamiento sin necesidad de acceder a los datos
originales.

El presente documento describe el funcionamiento de Qsimov y los experimentos realizados para
evaluar su comportamiento en escenarios de aprendizaje continual, siguiendo la estructura
recomendada para trabajos de fin de grado de ingeniería en informática~\cite{ref2}.

\section{Motivación y Objetivos}

\subsection{Motivación}

La motivación del proyecto nace de una observación fundamental sobre la naturaleza del entrenamiento
de redes neuronales con activación ReLU. Una red feed-forward con ReLU define, implícitamente, una
partición del espacio de entrada en regiones lineales: dentro de cada región, la función que computa
la red es una transformación lineal afín~\cite{ref3}. Esa transformación queda completamente
determinada por el subconjunto de neuronas que están activas para las entradas de esa región, es
decir, por el <<camino activo>> de la muestra.

Esta propiedad tiene una consecuencia directa para el re-entrenamiento: si se fijan los parámetros
que determinan \textit{qué} caminos existen (los pesos del modelo base, aprendidos durante el
pre-entrenamiento), el problema de encontrar los pesos óptimos para esos caminos se reduce a
resolver un sistema de ecuaciones lineales. Y un sistema lineal puede ampliarse con nuevas
ecuaciones en cualquier momento, incorporando nuevas muestras o nuevas clases, sin que las
soluciones previas se destruyan, siempre que las ecuaciones antiguas permanezcan en el sistema.

Frente a las soluciones existentes basadas en penalizaciones (EWC~\cite{ref4}, SI~\cite{ref5}) o
en expansión de arquitectura (PNN~\cite{ref6}), Qsimov ofrece garantías exactas de no-olvido para
las muestras incluidas en el sistema, sin requerir hiperparámetros adicionales y sin aumentar el
tamaño de la red.

\subsection{Objetivos}

Los objetivos del proyecto son los siguientes:

El objetivo principal del proyecto es evaluar experimentalmente la librería Qsimov en escenarios de
aprendizaje continual, caracterizando su comportamiento frente al olvido catastrófico y
comparándolo con métodos de re-entrenamiento estándar.

Los objetivos secundarios son:

\begin{enumerate}
  \item Comprender el funcionamiento de los componentes principales de Qsimov
        (\texttt{PathSelector}, \texttt{QsimovLinearSystem}, \texttt{QsimovGradient}) y su API de
        uso.
  \item Diseñar y ejecutar experimentos de evaluación sobre tres \textit{benchmarks} de distinta
        escala y complejidad: MNIST, CIFAR-10 e ImageNet (100 clases).
  \item Comparar el rendimiento de Qsimov (variante de sistema lineal y variante de gradiente) con
        métodos de referencia: fine-tuning estándar y re-entrenamiento acumulado.
  \item Analizar el impacto del punto de corte (\texttt{initial\_layer}) sobre la precisión
        obtenida y el número de caminos generados.
  \item Evaluar la robustez al learning rate del módulo de gradiente frente al re-entrenamiento
        estándar.
  \item Identificar los escenarios en los que Qsimov ofrece ventajas claras y aquellos en los que
        presenta limitaciones.
\end{enumerate}
